\begin{tikzpicture}[
    >=Stealth,
    node distance=1.8cm and 1.5cm,
    every node/.style={
        circle,
        draw=cyan!20!black,
        thick,
        minimum size=10mm,
        inner sep=0pt,
        fill=gray!10
    },
    edge/.style={
        ->,
        thick
    },
    active/.style={
        fill=red!85,
        draw=cyan!20!black
    },
    pathline/.style={
        red,
        ultra thick
    }
]

% Layer 0
\node (s) at (0,0) {$S$};

% Layer 1
\node (a1) at (2,3) {$a_1$};
\node (a2) at (2,1.5) {$a_2$};
\node (a3) at (2,0) {$a_3$};
\node (a4) at (2,-1.5) {$a_4$};
\node (a5) at (2,-3) {$a_5$};

% Layer 2
\node (b1) at (4,4) {$b_1$};
\node (b2) at (4,2) {};
\node (b3) at (4,0) {};
\node (b4) at (4,-2) {};
\node (b5) at (4,-4) {};

% Layer 3
\node (c1) at (6,4) {$c_1$};
\node (c2) at (6,2) {};
\node (c3) at (6,0) {};
\node (c4) at (6,-2) {};
\node (c5) at (6,-4) {};

% Layer 4
\node (d1) at (8,2) {};
\node (d2) at (8,0) {};
\node (d3) at (8,-2) {};
\node (d4) at (8,-4) {};

% Layer 5
\node (e1) at (10,1) {};
\node (e2) at (10,-1) {};
\node (e3) at (10,-3.5) {};

% Layer 6
\node (f1) at (12,2) {};
\node (f2) at (12,0) {};
\node (f3) at (12,-2) {};
\node (f4) at (12,-4) {};

% Edges
\draw[edge] (s) -- (a1);
\draw[edge] (s) -- (a2);
\draw[edge] (s) -- (a3);
\draw[edge] (s) -- (a4);
\draw[edge] (s) -- (a5);

\draw[edge] (a1) -- (b1);
\draw[edge] (a2) -- (b2);
\draw[edge] (a3) -- (b3);
\draw[edge] (a4) -- (b4);
\draw[edge] (a5) -- (b5);

\draw[edge] (b1) -- (c1);
\draw[edge] (b2) -- (c2);
\draw[edge] (b3) -- (c3);
\draw[edge] (b4) -- (c4);
\draw[edge] (b5) -- (c5);

% \draw[edge] (c1) -- (d1);
\draw[edge] (c2) -- (d1);
\draw[edge] (c2) -- (d2);
\draw[edge] (c3) -- (d3);
\draw[edge] (c4) -- (d4);

\draw[edge] (d2) -- (e1);
\draw[edge] (d2) -- (e2);
\draw[edge] (d3) -- (e2);
\draw[edge] (d3) -- (e3);

\draw[edge] (e1) -- (f1);
\draw[edge] (e2) -- (f2);
\draw[edge] (e3) -- (f3);
\draw[edge] (e3) -- (f4);

% Highlighted propagation paths
% \draw[pathline]
% (s.center) --
% (a2.center) --
% (b2.center) --
% (c2.center) --
% (d1.center);

% \draw[pathline]
% (s.center) --
% (a2.center) --
% (b2.center) --
% (c2.center) --
% (d2.center);

% \draw[pathline]
% (s.center) --
% (a3.center) --
% (b3.center) --
% (c3.center) --
% (d3.center);

\end{tikzpicture}